\documentclass[11pt]{article}
\usepackage[margin=1in]{geometry}
\usepackage{amsmath}
\usepackage{amssymb}
\usepackage{booktabs}
\usepackage{array}

\title{CSE 753 Advanced Reinforcement Learning\\
Study Notes: Imitation Learning and Offline RL\\
\large (Lecture 08, page by page)}
\author{}
\date{}

\begin{document}
\maketitle
\sloppy

\section{Slide 2: What Is Imitation Learning?}

\paragraph{What is the concept?}
Imitation Learning (IL) is a way of teaching a computer program (a \emph{policy}) to
perform a task by showing it recordings of an expert doing that task well --- instead of
telling the program a reward or a rule for what ``good'' means.

\paragraph{Explanation.}
Imagine you want to teach a robot to drive. You do not write down a mathematical rule
for ``good driving.'' Instead, you record hours of a human driving: what the car's
sensors saw (the \emph{state}) and what the driver did with the steering wheel and pedals
(the \emph{action}) at every instant. You then ask a computer to look at all of those
recordings and learn a rule that maps ``what the car sees'' to ``what to do,'' by copying
what the human did.

\paragraph{The equation.}
$$\mathcal{D} = \{(s_1, a_1, \ldots, s_T)\}$$

\begin{center}
\begin{tabular}{@{}p{3.2cm}p{10.3cm}@{}}
\toprule
Symbol & What it means \\
\midrule
$s_t$ & The state (what the world / sensors look like) at time step $t$. \\
$a_t$ & The action the expert took at time step $t$. \\
$T$ & The number of time steps in one recorded trajectory (one ``trip''). \\
$\mathcal{D}$ & The whole dataset: the full collection of recorded trajectories. \\
$\pi_{\text{expert}}$ & The unknown rule the expert was actually using to choose actions. \\
$\pi_\theta$ & The policy we want to train, with adjustable numbers (parameters) $\theta$. \\
\bottomrule
\end{tabular}
\end{center}

\paragraph{What the equation is saying.}
$\mathcal{D}$ is simply a list of recorded trips. Each trip is a sequence of
``what did the world look like, what did the expert do'' pairs. We never get to see the
expert's internal rule $\pi_{\text{expert}}$ directly --- only the footprints it left
behind in the data.

\paragraph{Aspects of this concept.}
\begin{itemize}
\item \textbf{Problem it solves:} it lets us build a working policy for tasks where
writing down a reward function is hard (e.g., what number captures ``good driving''?),
as long as we can get someone to demonstrate the task.
\item \textbf{Problem it cannot solve:} it only tells us what action the expert chose,
never how good that action was compared to alternatives; a policy trained this way can
never (in principle) beat the expert it copied.
\item \textbf{Problems it gives rise to:} if we are not careful about the learning
method, we run into the ``mean-collapse'' problem (Slide 4) and the ``compounding
error'' problem (Slide 12) --- both covered later.
\item \textbf{When to use it:} when demonstrations are available or cheap to collect,
and a good reward function is hard to write down.
\item \textbf{When not to use it:} when there is no expert to copy, or when the goal is
to do \emph{better} than any available expert --- that calls for reinforcement learning
(RL) instead, which is what the rest of this lecture builds toward.
\item \textbf{What if this were done differently?} If, instead of only $(s,a)$ pairs, the
dataset also included a reward $r$ at every step, we could use offline RL (Slide 15
onward) and potentially outperform the expert, rather than just copy it.
\end{itemize}

\section{Slide 3: Imitation Learning --- Version 0 (Behavior Cloning, Deterministic)}

\paragraph{What is the concept?}
The simplest possible way to do imitation learning: treat it as an ordinary supervised
regression problem. For every recorded state, make the policy's predicted action match
the expert's recorded action as closely as possible. This is called \textbf{Behavior
Cloning (BC)}.

\paragraph{Explanation.}
Think of it like a fill-in-the-blank exercise: you show the network a state, it guesses
an action, and you tell it exactly how wrong the guess was compared to what the expert
actually did. You repeat this over the whole dataset, nudging the network's numbers a
little bit each time so its guesses get closer to the expert's actions.

\paragraph{The equation.}
$$\min_\theta \ \frac{1}{|\mathcal{D}|}\sum_{(s,a)\in\mathcal{D}} \|a - \hat a\|^2,
\qquad \hat a = \pi_\theta(s)$$

\begin{center}
\begin{tabular}{@{}p{3.2cm}p{10.3cm}@{}}
\toprule
Symbol & What it means \\
\midrule
$\theta$ & The adjustable numbers (parameters) inside the policy network. \\
$|\mathcal{D}|$ & The total number of $(s,a)$ pairs in the dataset. \\
$a$ & The action the expert actually took in state $s$. \\
$\hat a$ & The action our policy predicts for that same state $s$. \\
$\pi_\theta(s)$ & The policy: a function that takes one state in and outputs one action. \\
$\|\cdot\|^2$ & Squared distance between two action vectors (how far apart they are). \\
$\min_\theta$ & We search for the parameters $\theta$ that make the total error smallest. \\
\bottomrule
\end{tabular}
\end{center}

\paragraph{What the equation is saying.}
For every recorded state, measure how far the policy's guess is from what the expert
did, square that distance (so being very wrong is punished a lot more than being a
little wrong), and average this over the whole dataset. Then adjust $\theta$ to make
this average as small as possible.

\paragraph{Aspects of this concept.}
\begin{itemize}
\item \textbf{Advantage:} extremely simple --- it is exactly ordinary supervised
regression, so all the standard machine-learning tools apply directly; fast and easy to
train.
\item \textbf{Disadvantage:} it assumes there is basically one ``correct'' action per
state. It cannot represent a situation where two very different actions are both
perfectly valid.
\item \textbf{Problem it solves:} turns policy learning into a completely standard,
well-understood regression problem --- no reward, no exploration, no RL loop needed.
\item \textbf{Problem it cannot solve:} if the expert data is \emph{multimodal} (the
expert sometimes did one thing and sometimes a very different thing in a similar
situation), this method collapses both choices into their average, which can be a
useless or dangerous action (see Slide 4).
\item \textbf{What problem it gives rise to:} the mean-collapse failure, which motivates
version 1 (Slide 6) with more expressive policies.
\item \textbf{When to use it:} when the expert's behavior is essentially
\emph{unimodal} --- one sensible action per state, with only small variation.
\item \textbf{When not to use it:} when the expert's behavior is multimodal, i.e.\ there
are multiple, very different, equally good actions in the same state.
\item \textbf{What if this were done differently?} If we swapped the squared-error loss
for a maximum-likelihood loss with an expressive distribution, we would get version 1
(Slide 6), which is able to represent more than one good action per state.
\end{itemize}

\section{Slide 4: What Can Go Wrong? The Mean-Collapse Problem}

\paragraph{What is the concept?}
This slide names and explains the specific failure of version-0 behavior cloning: when
the expert's actions are multimodal (more than one distinct, valid strategy), squared-error
regression learns the \emph{average} of those strategies, which is often an action that
is unsafe and was never actually taken by anyone.

\paragraph{Explanation.}
Picture a highway scene: some human drivers, seeing a slow truck ahead, merge left; other
drivers stay in their lane. Both are perfectly good choices. If you train a
squared-error policy on this mixed data, it will not learn ``sometimes merge, sometimes
stay.'' It will learn one single number that sits roughly \emph{between} those two
choices --- which, on a real road, means steering the car into the gap between the lanes,
straight toward the truck. That compromise action is worse than either of the two things
any actual human did.

\paragraph{Mechanism (why this happens).}
For a fixed input, squared-error loss is minimized exactly when the prediction equals
the \emph{mean} of the target values for that input. This is a basic fact about squared
error, not something specific to neural networks. If the true action distribution
$p(a\mid s)$ has two separate ``bumps'' (modes), the mean falls in the valley between
them --- a point with almost no real probability under the data.

\paragraph{Worked numerical example.}
Suppose $75\%$ of recorded examples for a given state are ``stay straight'' ($a\approx
0^\circ$) and $25\%$ are ``merge left'' ($a\approx -2^\circ$ steering angle). The learned
(mean) action is
$$\mathbb{E}[a\mid s] = 0.75\times 0 + 0.25\times(-2) = -0.5^\circ,$$
which matches the dashed line on the slide. This $-0.5^\circ$ action was never actually
demonstrated by anyone; it straddles both lanes.

\paragraph{Aspects of this concept.}
\begin{itemize}
\item \textbf{Name of the failure:} mean-seeking behavior (also called mode-averaging or
mean-collapse).
\item \textbf{Why it is dangerous, specifically:} the learned action is not a safe
compromise between two valid strategies --- it is a third action that lies in the gap
between them and that no expert ever took. On a highway, it drives into the obstacle; in
a surgical setting, it could mean cutting exactly where neither valid surgical approach
cuts.
\item \textbf{Problem it gives rise to:} the need for expressive (multimodal) policy
classes, covered starting at Slide 5 and built in Slides 6--11.
\item \textbf{What if this were done differently?} If the policy's output distribution
could represent two separate bumps instead of one number, it could put its probability
mass on both ``merge'' and ``stay,'' rather than averaging them. This is exactly what
Slides 5--11 build.
\end{itemize}

\section{Slide 5: Neural Networks as Distribution Parameterizations}

\paragraph{What is the concept?}
A neural network only ever outputs numbers; how those numbers are turned into an action
depends on what kind of probability distribution you decided to wrap around the output.
The expressive power of the \emph{network} (how complicated a function it can compute)
is a completely different question from the expressive power of the \emph{distribution
family} you chose for its output.

\paragraph{Explanation.}
For a video-game character choosing among a small, fixed list of buttons (up, down, A,
B), the network can output one probability per button. This is called a
\emph{categorical} distribution, and it can represent \emph{any} pattern of preferences
over those buttons --- including ``sometimes A, sometimes B, rarely up or down.'' That
is maximally expressive. But for a continuous action like a steering angle, the slide's
simple design has the network output just two numbers, $\mu$ (the center) and $\sigma$
(the spread), of a single bell curve (Gaussian). A single bell curve can only ever have
\emph{one} peak. No matter how large or clever the network is, if you force its output
through a single Gaussian, it can never represent ``sometimes steer left, sometimes stay
straight'' as two separate peaks --- it will always squash them into one peak in the
middle, which is exactly the mean-collapse problem from Slide 4.

\paragraph{Aspects of this concept.}
\begin{itemize}
\item \textbf{Advantage of the categorical (discrete) case:} naturally and fully
multimodal --- no extra machinery needed.
\item \textbf{Disadvantage of the single-Gaussian (continuous) case:} unimodal by
construction, regardless of how big or well-trained the underlying network is.
\item \textbf{Problem it explains:} \emph{why} Slide 3/4's version-0 policy fails for
continuous multimodal experts --- it is not a training failure, it is a
representational limitation baked into the choice of a single Gaussian.
\item \textbf{What problem it gives rise to:} the need for continuous-action
distribution families that \emph{can} have more than one peak --- Gaussian Mixture
Models, autoregressive models, and diffusion/flow policies (Slides 7--11).
\item \textbf{When to use a single Gaussian:} only when you are confident the expert
behavior really is unimodal for continuous actions.
\item \textbf{What if this were done differently?} If the continuous-action output were
instead a mixture of several Gaussians (Slide 7), the same network architecture idea
(numbers in, numbers out) becomes just as expressive as the categorical case.
\end{itemize}

\section{Slide 6: Imitation Learning --- Version 1 (Expressive Policies)}

\paragraph{What is the concept?}
Instead of regressing to a single number (version 0), train the policy to output a full
probability distribution over actions, and make the expert's actual action as
\emph{likely as possible} under that distribution. This is \textbf{maximum likelihood
estimation (MLE)}, and it opens the door to multimodal policies.

\paragraph{Explanation.}
Rather than saying ``get as close as possible to this exact number,'' we now say ``make
this exact action look as probable as possible under whatever shape of distribution you
are outputting.'' If the distribution family is expressive enough (can have more than
one peak), the network can learn to put a peak over \emph{each} of the expert's distinct
strategies, instead of averaging them together.

\paragraph{The equation.}
$$\min_\theta \ -\mathbb{E}_{(s,a)\in\mathcal{D}}\left[\log \pi_\theta(a\mid s)\right]$$

\begin{center}
\begin{tabular}{@{}p{3.2cm}p{10.3cm}@{}}
\toprule
Symbol & What it means \\
\midrule
$\theta$ & Adjustable parameters of the policy network. \\
$\mathbb{E}_{(s,a)\in\mathcal{D}}$ & Average over every $(s,a)$ pair in the dataset. \\
$\pi_\theta(a\mid s)$ & The probability (or probability density) the policy assigns to
action $a$, given state $s$; can be any shape, not just a single bell curve. \\
$\log \pi_\theta(a\mid s)$ & The log-probability of the expert's actual action; higher
means the policy thinks that action was more likely. \\
$-\mathbb{E}[\ldots]$ & We minimize the \emph{negative} average log-probability, which is
the same as maximizing the average log-probability. \\
\bottomrule
\end{tabular}
\end{center}

\paragraph{What the equation is saying.}
Look at every expert action in the dataset. Push the policy's parameters so that the
policy assigns \emph{high probability} to exactly the actions the expert took, on
average across the whole dataset.

\paragraph{Aspects of this concept.}
\begin{itemize}
\item \textbf{Advantage:} more general than version 0; if paired with an expressive
distribution family, it can represent multiple distinct valid strategies instead of
averaging them.
\item \textbf{Disadvantage:} by itself, this objective does not automatically fix
mean-collapse --- everything depends on how expressive $\pi_\theta(\cdot\mid s)$ is. If
$\pi_\theta$ is still a single Gaussian, maximizing likelihood gives essentially the same
mean-collapse result as version 0.
\item \textbf{Problem it solves:} provides the general recipe (maximum likelihood) that
Slides 7--11 plug different expressive distribution families into.
\item \textbf{Problem it does not solve:} it does not address compounding error /
distribution shift (Slide 12), which is a separate issue about \emph{which states} the
policy will see, not about how expressive its action distribution is.
\item \textbf{When to use it:} whenever the expert's behavior might be multimodal, i.e.\
almost always for real-world tasks with more than one reasonable strategy.
\item \textbf{What if this were done differently?} If we kept this exact objective but
plugged in a single Gaussian for $\pi_\theta$, we would be back to (a probabilistic
version of) version 0 --- proof that the objective alone is not what fixes mean-collapse;
the distribution family matters just as much.
\end{itemize}

\section{Slide 7: Gaussian Mixture Models (GMMs)}

\paragraph{What is the concept?}
A Gaussian Mixture Model represents the action distribution as a weighted combination of
several separate bell curves (Gaussians), each one representing a different possible
``strategy.''

\paragraph{Explanation.}
Imagine several bell-shaped hills, each centered on a different action (a different
strategy the expert might use), each with its own width and its own ``how common is
this strategy'' weight. Add all these hills together (after scaling each by its weight)
and you get one bumpy probability curve that can have several separate peaks --- one per
strategy --- instead of being forced into a single hill.

\paragraph{The equations.}
$$\{w_k(s),\ \mu_k(s),\ \Sigma_k(s)\}_{k=1}^K \quad(1)\qquad\qquad
p(a_i\mid s_i) = \sum_{k=1}^{K} w_k(s_i)\,\mathcal{N}\!\left(a_i \mid \mu_k(s_i), \Sigma_k(s_i)\right)\quad(2)$$
$$\max_\theta \sum_{i=1}^N \log p(a_i\mid s_i) \quad(3)\qquad\qquad
\mathcal{L}(\theta) = -\sum_{i=1}^N \log\left[\sum_{k=1}^K w_k(s_i)\,\mathcal{N}(a_i\mid\mu_k(s_i),\Sigma_k(s_i))\right]\quad(4)$$

\begin{center}
\begin{tabular}{@{}p{3.2cm}p{10.3cm}@{}}
\toprule
Symbol & What it means \\
\midrule
$K$ & The number of Gaussian ``strategies'' (mixture components) the network can use;
chosen ahead of time. \\
$w_k(s)$ & How likely strategy $k$ is, given state $s$ (all $w_k$ for a fixed $s$ add up
to 1). \\
$\mu_k(s)$ & The center (typical action) of strategy $k$, given state $s$. \\
$\Sigma_k(s)$ & The spread/uncertainty (covariance) of strategy $k$, given state $s$. \\
$\mathcal{N}(a\mid\mu,\Sigma)$ & The height of a Gaussian bell curve with center $\mu$
and spread $\Sigma$, evaluated at action $a$. \\
$p(a_i\mid s_i)$ & The overall probability density of action $a_i$: the weighted sum of
all $K$ strategies' bell curves. \\
$N$ & The number of $(s,a)$ training examples. \\
$\mathcal{L}(\theta)$ & The negative log-likelihood loss actually minimized during
training (equivalent to maximizing eq.\ (3)). \\
\bottomrule
\end{tabular}
\end{center}

\paragraph{What the equations are saying.}
Equation (1): for every state, the network outputs $K$ separate ``strategies,'' each
with its own weight, center, and spread. Equation (2): the overall probability of an
action is a weighted blend of all $K$ bell curves. Equations (3)--(4): train the network
so that, averaged over the dataset, the expert's actual actions get high probability
under this blended, multi-peaked curve.

\paragraph{Worked numerical example (checking the notation).}
Suppose two strategies have means $-30^\circ$ and $+30^\circ$, each with standard
deviation $\sigma = 2$ (so $\Sigma = \sigma^2 = 4$, since the second number inside
$\mathcal{N}(\cdot\,;\text{mean},\text{variance})$ is the \emph{variance}, not the
standard deviation). If we ask for the density at $a=0^\circ$ (right in between the two
peaks), each Gaussian's density is essentially zero there, because $0^\circ$ is $15$
standard deviations away from either mean ($30/2 = 15$). So $p(a=0\mid s)\approx 0.000$
--- $a=0$ sits in the trough between the modes, exactly the unsafe ``average'' action
that Slide 4 warned about. A GMM correctly assigns it almost no probability, whereas a
single Gaussian (Slide 5) would have centered its one peak right there.

\paragraph{Aspects of this concept.}
\begin{itemize}
\item \textbf{Advantage:} can represent several genuinely distinct action strategies at
once; the density has a clean closed-form formula, is easy to sample from, and is easy
to differentiate for training.
\item \textbf{Disadvantage:} the number of strategies $K$ must be fixed in advance
(cannot grow on its own), and it scales poorly to actions with many dimensions (the
covariance matrices get large and hard to estimate).
\item \textbf{Problem it solves:} directly fixes the mean-collapse problem from Slide 4
by letting the output have multiple peaks.
\item \textbf{Problem it cannot solve:} it does not, by itself, model rich dependencies
between different dimensions of a single action vector unless the covariance $\Sigma_k$
is allowed to be a full (not just diagonal) matrix, which is expensive; it also does not
address compounding error.
\item \textbf{When to use it:} when you have a good guess of how many distinct
strategies exist (a small, known $K$) and actions are low-to-moderate dimensional.
\item \textbf{When not to use it:} for very high-dimensional, richly structured actions
where the individual dimensions strongly depend on one another --- an autoregressive
model (Slide 8) is often a better fit there.
\item \textbf{What if this were done differently?} If $K=1$, a GMM is exactly a single
Gaussian --- back to the non-expressive case of Slide 5. If $K$ is set far larger than
the number of real strategies, the model can overfit or become unstable to train.
\end{itemize}

\section{Slide 8: Multimodal Actions --- The Chain-Rule (Autoregressive) Idea}

\paragraph{What is the concept?}
Instead of representing a multi-dimensional action all at once (like a GMM does), break
the action into its individual dimensions and predict them one at a time, each new
dimension conditioned on all the dimensions already chosen. This uses the
\textbf{chain rule of probability}, which is always mathematically exact.

\paragraph{Explanation.}
Suppose an action has several parts (e.g., steering angle \emph{and} speed). Instead of
guessing both at once, first guess the steering angle. Then, \emph{knowing} what
steering angle was just chosen, guess the speed (a sharp turn probably means a slower
speed, for example). This lets the model capture how the different parts of an action
relate to each other, which a single joint distribution might miss if it treats the
dimensions independently.

\paragraph{The equations.}
$$a_i = (a_{i,1}, a_{i,2}, \ldots, a_{i,D}) \quad(5) \qquad\qquad
p_\theta(a_i\mid s_i) = \prod_{d=1}^{D} p_\theta\!\left(a_{i,d}\mid s_i, a_{i,1:d-1}\right) \quad(6)$$
$$\max_\theta \sum_{i=1}^N \log p_\theta(a_i\mid s_i) \quad(7) \qquad\qquad
\mathcal{L}(\theta) = -\sum_{i=1}^N \sum_{d=1}^D \log p_\theta\!\left(a_{i,d}\mid s_i, a_{i,1:d-1}\right)\quad(8)$$

\begin{center}
\begin{tabular}{@{}p{3.2cm}p{10.3cm}@{}}
\toprule
Symbol & What it means \\
\midrule
$D$ & The number of separate dimensions (parts) in one action vector. \\
$a_{i,d}$ & The $d$-th dimension of the $i$-th training action. \\
$a_{i,1:d-1}$ & All the dimensions of action $i$ chosen \emph{before} dimension $d$. \\
$p_\theta(a_{i,d}\mid s_i, a_{i,1:d-1})$ & The probability of dimension $d$, given the
state and every dimension already decided. \\
$\prod_{d=1}^D$ & Multiply the per-dimension probabilities together (chain rule). \\
\bottomrule
\end{tabular}
\end{center}

\paragraph{What the equations are saying.}
Equation (5) just says an action can have several parts. Equation (6) says the
probability of the whole action equals the probability of the first part, times the
probability of the second part given the first, times the probability of the third part
given the first two, and so on --- this is always exactly true for any joint
distribution, by the definition of conditional probability. Equations (7)--(8) say: train
by maximizing (equivalently, minimizing the negative of) the total log-probability of the
expert's actual action, summed over every dimension and every training example.

\paragraph{Aspects of this concept.}
\begin{itemize}
\item \textbf{Advantage:} the factorization in equation (6) is \emph{exact} --- no
information is lost --- and it explicitly captures dependencies between different
dimensions of the same action, which a simple GMM with independent dimensions cannot.
\item \textbf{Disadvantage:} actions must be generated one dimension at a time in
sequence, which is slower at execution time than predicting everything in one shot; an
error in an early dimension can affect every later dimension of the \emph{same} action.
\item \textbf{Problem it solves:} models rich correlations \emph{within} a single
action (e.g., steering and speed changing together), not just multiple separate
strategies.
\item \textbf{When to use it:} high-dimensional, structured actions where the
dimensions clearly interact (e.g., the joint angles of a robot arm).
\item \textbf{What if this were done differently?} If each dimension were instead
predicted independently of the others (no conditioning on $a_{i,1:d-1}$), the model
could not represent any relationship between dimensions --- e.g., it could never learn
that ``sharp turn implies slower speed.''
\end{itemize}

\section{Slide 9: Autoregressive Models --- Generating the Action}

\paragraph{What is the concept?}
This slide shows how an autoregressive policy actually \emph{produces} an action once
trained: dimension by dimension, sampling each one from its own conditional
distribution.

\paragraph{Explanation.}
Once trained, using the policy is like a chain of small decisions: pick dimension 1;
now that dimension 1 is fixed, pick dimension 2 knowing dimension 1; keep going until
every dimension of the action has been chosen.

\paragraph{The equations.}
$$\pi_\theta(a\mid s) = \prod_{d=1}^D \pi_\theta(a_d \mid s, a_{<d}) \quad(9)
\qquad\qquad
s \to a_1 \to a_2 \to \cdots \to a_D \quad(10)$$

\begin{center}
\begin{tabular}{@{}p{3.2cm}p{10.3cm}@{}}
\toprule
Symbol & What it means \\
\midrule
$a_{<d}$ & Shorthand for $(a_1,\ldots,a_{d-1})$: everything chosen before dimension $d$. \\
$a_1 \sim \pi_\theta(a_1\mid s)$ & Sample the first dimension from its distribution,
given only the state. \\
$a_2 \sim \pi_\theta(a_2\mid s,a_1)$ & Sample the second dimension, now also conditioned
on the first dimension just chosen. \\
$s \to a_1 \to \cdots \to a_D$ & The generation order: state first, then each action
dimension in turn. \\
\bottomrule
\end{tabular}
\end{center}

\paragraph{What the equations are saying.}
Equation (9) is identical in spirit to equation (6) on the previous slide --- it is the
same chain-rule factorization, just written for execution rather than training.
Equation (10) shows this as a left-to-right generation chain: state in, then one
dimension of the action out at a time, each new dimension informed by all the previous
ones.

\paragraph{Aspects of this concept.}
This slide is a direct continuation of Slide 8, so the advantages, disadvantages, and
use cases are the same. The one new point worth remembering: this factorization
explicitly models dependencies between action dimensions, and the \textbf{Decision
Transformer} (next slide) is a concrete, well-known architecture built on exactly this
autoregressive idea.

\section{Slide 10: Decision Transformer}

\paragraph{What is the concept?}
A concrete architecture that turns imitation learning into a sequence-modeling problem
(the same kind of problem language models solve), by feeding a Transformer a sequence
of (desired future return, state, action) triples and asking it to predict the next
action.

\paragraph{Explanation.}
Instead of thinking in terms of states and rewards the usual RL way, treat the whole
problem like predicting the next word in a sentence. The ``sentence'' here is a sequence
of tokens: ``how much reward do I still want to get,'' ``what do I currently see,''
``what did I just do,'' repeated over time. A \emph{causal} Transformer (one that can
only look at the past, never the future) reads this sequence and predicts the next
action.

\begin{center}
\begin{tabular}{@{}p{3.2cm}p{10.3cm}@{}}
\toprule
Symbol & What it means \\
\midrule
$\hat R_t$ & The return-to-go: how much total reward we would like to still receive from
this point onward. \\
$s_t$ & The state at time $t$. \\
$a_t$ & The action at time $t$. \\
Causal transformer & A Transformer that, when predicting position $t$, is only allowed
to look at positions $1,\ldots,t$, never anything later. \\
Linear decoder & The final small layer that turns the Transformer's internal numbers
into an actual action. \\
\bottomrule
\end{tabular}
\end{center}

\paragraph{Aspects of this concept.}
\begin{itemize}
\item \textbf{Advantage:} reuses powerful, well-understood sequence models
(Transformers); at test time, you can ask for a high desired return and the model will
try to act like a higher-performing agent, entirely through ordinary supervised training
(no bootstrapped value function needed).
\item \textbf{Disadvantage:} needs a reasonably long context window and has the usual
quadratic-in-length cost of attention; because it only imitates sequences that were
actually observed, asking for an unrealistically high return can make it behave poorly
(it has nothing real to imitate at that return level).
\item \textbf{Problem it solves:} avoids the instability that comes from bootstrapped
value learning (TD updates), replacing it with plain supervised sequence prediction.
\item \textbf{Problem it does not solve:} because it does not perform explicit value
backups, it generally cannot \emph{stitch} together good pieces of different
trajectories into a new, better trajectory the way TD-based offline RL can (see Slide
16).
\item \textbf{When to use it:} when you have a large offline dataset of trajectories
with a range of returns and want a simple, stable, supervised training recipe.
\item \textbf{What if this were done differently?} If the Transformer were
\emph{non-causal} (allowed to look at future tokens), it would be cheating by peeking at
information it will not have during real-time control --- the causal mask is essential
for the model to be usable at actual decision time.
\end{itemize}

\section{Slide 11: Diffusion Policy}

\paragraph{What is the concept?}
A way to turn random noise into a realistic action by learning a ``velocity field'' that
gradually pushes noisy points toward the shape of the real action distribution.

\paragraph{Note on the slide.}
This slide is titled ``Diffusion Policy,'' but the actual procedure shown --- a straight
linear interpolation between noise and data, and a single learned velocity field
integrated with simple Euler steps --- is what is technically called
\textbf{flow matching} (or rectified flow), not classic diffusion (DDPM). Classic
diffusion gradually corrupts data with many small noise steps and learns to reverse that
corruption stochastically. Flow matching, shown here, learns one direct ``straight-line''
velocity field instead. For the exam, answer using exactly the procedure written on the
slide; if you want to show extra precision, you can note in one sentence that the
mechanism is really a learned vector field (flow matching), not iterative denoising.

\paragraph{Explanation.}
Start with a completely random point (pure noise). Learn a ``which direction should I
walk'' rule (the velocity field) that, if you take many small steps following it, will
carry that random point all the way to a realistic action. Different random starting
points end up at different final actions --- so this approach is naturally multimodal:
different noise samples can land near different expert strategies.

\paragraph{The steps (training and inference).}

\begin{center}
\begin{tabular}{@{}p{3.2cm}p{10.3cm}@{}}
\toprule
Symbol & What it means \\
\midrule
$x_1 \sim \mathcal{D}$ & A real action sampled from the dataset. \\
$x_0 \sim \mathcal{N}(0,1)$ & A random noise sample. \\
$t \sim p(t)$ & A random ``progress fraction'' between 0 (all noise) and 1 (all data). \\
$x_t = t x_1 + (1-t) x_0$ & A blend of noise and data: the point that is $t$ fraction of
the way from noise to data. \\
$v_\theta(x_t,t)$ & The network's prediction of ``which direction to move'' at point
$x_t$ and progress $t$ (conditioned on the state $s_t$ for imitation learning). \\
$\delta$ & The step size used when numerically integrating the velocity field. \\
\bottomrule
\end{tabular}
\end{center}

\textbf{Training:} sample $x_1\sim\mathcal{D}$ and $x_0\sim\mathcal{N}(0,1)$; sample $t$;
interpolate $x_t = tx_1+(1-t)x_0$; minimize $\|v_\theta(x_t,t)-(x_1-x_0)\|^2$ (i.e., teach
the network to predict the straight-line direction from noise to data).

\textbf{Inference:} start from noise $x_0\sim\mathcal{N}(0,1)$; repeatedly step
$x_{t+\delta} = x_t + v_\theta(x_t,t)\delta$ for $t=0,\delta,2\delta,\ldots,1-\delta$;
return $x_1$ as the final action.

\paragraph{Aspects of this concept.}
\begin{itemize}
\item \textbf{Advantage:} highly expressive --- can represent very complex,
high-dimensional, arbitrarily-shaped multimodal action distributions, more flexible than
a fixed-$K$ GMM.
\item \textbf{Disadvantage:} generating one action requires several iterative steps
(slower at test time than a single forward pass through a GMM or autoregressive model);
more moving parts to implement and tune (step size, number of steps).
\item \textbf{Problem it solves:} represents very rich, complex continuous multimodal
distributions that a fixed-size GMM or a simple autoregressive model may struggle to
capture.
\item \textbf{Problem it does not solve:} like every method in Slides 6--11, it only
improves \emph{how expressive} the action distribution is; it does nothing about
compounding error / distribution shift (Slide 12), which is a separate, data-coverage
problem.
\item \textbf{When to use it:} complex, high-dimensional continuous control (e.g.,
robot manipulation) where test-time compute for several iterative steps is affordable.
\item \textbf{When not to use it:} latency-critical, real-time control where only a
single fast forward pass is acceptable.
\item \textbf{What if this were done differently?} Using a larger step size $\delta$
(fewer integration steps) makes generation faster but less accurate (bigger
discretization error) --- a direct speed-versus-accuracy trade-off.
\end{itemize}

\section{Slide 12: Problems with Imitation Learning --- Compounding Error}

\paragraph{What is the concept?}
Even with a perfectly expressive policy (Slides 6--11), imitation learning has a second,
completely different failure mode: small mistakes can snowball because the policy's own
actions change what state it sees next.

\paragraph{Explanation.}
In ordinary supervised learning (like recognizing pictures of cats), each example is
independent --- one wrong guess does not change which picture you see next. But in
imitation learning, the policy's action actually changes the world: a slightly wrong
steering command puts the car in a slightly different position than any human driver was
ever recorded in. The policy was never trained on that slightly-off position, so it is
more likely to make another, possibly bigger, mistake there --- pushing the car even
further from anywhere in the training data. This keeps happening, and small errors
compound into large ones.

\paragraph{Mechanism, step by step (this is the full failure chain).}
\begin{enumerate}
\item The predicted action affects the next state; states are not independent of the
policy, unlike ordinary supervised learning.
\item A small prediction error nudges the agent slightly off the states the expert ever
demonstrated.
\item The policy was never trained on this new, off-distribution state, so its
prediction there is unreliable and tends to be more wrong.
\item This repeats at every subsequent step, and errors compound over the length of the
trajectory --- growing much faster than they would if each mistake were independent.
\end{enumerate}

\paragraph{Aspects of this concept.}
\begin{itemize}
\item \textbf{Name of the failure:} compounding error, also called distribution shift.
\item \textbf{Two remedies shown on the slide:}
\begin{enumerate}
\item \emph{Collect a lot of demonstration data and hope for the best.} Advantage:
requires no algorithm change, very simple. Disadvantage: does not really fix the
underlying issue --- it only reduces the \emph{chance} of drifting off-distribution; you
can never demonstrate every possible state in advance, especially in a large or
continuous state space.
\item \emph{Collect corrective behavior data} --- this leads directly to DAgger, the
next slide.
\end{enumerate}
\item \textbf{Problem it gives rise to:} the DAgger algorithm (Slide 13), and, more
broadly, the entire motivation for offline RL, which uses reward information to be more
robust to imperfect data coverage.
\end{itemize}

\section{Slide 13: DAgger --- Dataset Aggregation}

\paragraph{What is the concept?}
DAgger directly fixes compounding error by changing \emph{where} the training data comes
from: instead of only collecting states the expert visited, it collects states the
\emph{learner itself} visits, and labels those states with what the expert would have
done there.

\paragraph{Explanation.}
Run your current (still imperfect) policy in the real environment and see exactly where
it goes. At each of those states, go ask the real expert, ``what would you have done
here?'' Add these new (state, expert-action) pairs to your dataset, and retrain on
everything collected so far. Repeat this loop. Over time, the training data starts to
match exactly the states the policy will actually encounter, so there is no more gap
between training and deployment.

\paragraph{The equation.}
$$\mathcal{D} \leftarrow \mathcal{D} \cup \{(s', a^*)\}, \qquad s' \sim \pi_\theta,\quad
a^* \sim \pi_{\text{expert}}(\cdot\mid s')$$

\begin{center}
\begin{tabular}{@{}p{3.2cm}p{10.3cm}@{}}
\toprule
Symbol & What it means \\
\midrule
$s' \sim \pi_\theta$ & A state actually visited by rolling out the current learned
policy. \\
$a^* \sim \pi_{\text{expert}}(\cdot\mid s')$ & The action the real expert would take at
that visited state, queried on demand. \\
$\mathcal{D} \leftarrow \mathcal{D}\cup\{(s',a^*)\}$ & Add this new labeled pair to the
growing dataset (``aggregate''). \\
\bottomrule
\end{tabular}
\end{center}

\paragraph{Full loop.}
Roll out $\pi_\theta$ to get visited states $\to$ query the expert at those states $\to$
aggregate the new pairs into $\mathcal{D}$ $\to$ retrain $\min_\theta \mathcal{L}(\pi_\theta,\mathcal{D})$
$\to$ repeat.

\paragraph{Aspects of this concept.}
\begin{itemize}
\item \textbf{Advantage:} directly closes the mismatch between training data and
deployment data that caused compounding error --- the policy is trained exactly on the
states it will actually visit.
\item \textbf{Disadvantage:} requires an \emph{interactive} expert available throughout
training, one that can be asked to label any state on demand. This is often infeasible
--- a surgeon cannot realistically sit and label every single state a surgical robot
visits during training. The rollout-query-aggregate-retrain loop is also slow and
expensive to run repeatedly.
\item \textbf{Problem it solves:} eliminates distribution-shift-driven compounding
error at its root.
\item \textbf{Problem it cannot solve:} it does not help when no interactive expert is
available, or when the ``expert'' has no good answer for a truly bizarre state the
learner has drifted into.
\item \textbf{When to use it:} when a cheap, available interactive expert exists (e.g.,
a scripted expert in simulation, or a supervisor watching a live feed).
\item \textbf{When not to use it:} when only a fixed batch of past demonstrations exists
and no one can be asked new questions (e.g., historical human driving logs) --- offline
methods (Slides 14 onward) are the fallback there.
\item \textbf{What if this were done differently?} If step 2 used the \emph{old}
dataset's nearest recorded action instead of asking the true expert fresh, the new label
would not actually reflect what should be done at that specific new state --- the whole
point of DAgger is that the correction has to come from the real expert reacting live to
the state actually visited.
\end{itemize}

\section{Slide 14: Why Offline RL?}

\paragraph{What is the concept?}
This slide contrasts the usual online RL loop (collect data, update policy, repeat)
with offline RL, where the policy is trained once on a fixed, already-collected dataset
and never interacts with the real environment during training.

\paragraph{Explanation.}
Online RL is like learning to ride a bike by actually riding it, falling, and adjusting
--- over and over, in real time. Offline RL is like being handed a giant pile of videos
of other people riding bikes (some skilled, some not) and being asked to learn to ride
well just by studying those videos, without ever getting on a bike yourself during
training.

\begin{center}
\begin{tabular}{@{}p{4.5cm}p{9cm}@{}}
\toprule
Online RL & Offline RL \\
\midrule
Collect data, then update the policy on the latest data (or all data so far), and
repeat. & Given one static, already-collected dataset; train the policy purely on that
provided dataset. \\
\bottomrule
\end{tabular}
\end{center}

\paragraph{Aspects of this concept.}
\begin{itemize}
\item \textbf{Advantages of offline RL:} it can reuse data other people or systems
already collected; live data collection may be risky or unsafe (a self-driving car or a
surgical robot exploring randomly in the real world is dangerous); data can be reused
across different experiments, projects, robots, or institutions; and offline pretraining
can be blended with a later stage of online fine-tuning.
\item \textbf{Disadvantage:} the fixed dataset may not cover every state or action well,
and it may not contain optimal behavior --- this data-coverage gap is the seed of the
out-of-distribution (OOD) action problem discussed in Slides 17--18.
\item \textbf{Problem it solves:} lets us learn from safe, already-existing data instead
of expensive or dangerous live trials.
\item \textbf{Problem it does not solve:} it does not guarantee the data covers
everything the learned policy might need --- care is required so the policy does not
over-trust regions the data barely covers.
\item \textbf{When to use it:} safety-critical or expensive-to-explore domains where
logged data already exists.
\item \textbf{When not to use it:} domains where cheap, safe live exploration is
available --- online RL may simply be easier there.
\end{itemize}

\section{Slide 15: Offline RL --- Formal Setup}

\paragraph{What is the concept?}
This slide gives the precise mathematical definition of the offline RL problem: a
dataset generated by some unknown ``behavior policy,'' and a goal that is still the
usual RL goal --- maximize expected total reward under the policy we are training, not
under the policy that generated the data.

\paragraph{Explanation.}
We are handed a big table of (state, action, next state, reward) records, produced by
\emph{some} policy $\pi_\beta$ that we do not control and may not even fully know (it
could even be a mix of several different people's or systems' behavior). Even though we
only get to look at $\pi_\beta$'s data, our actual goal has not changed: we still want to
find the best possible policy $\pi_\theta$ to run in the real environment.

\paragraph{The equations.}
$$p_{\pi_\beta}(\tau) = p(s_1)\prod_{t=1}^{T}\pi_\beta(a_t\mid s_t)\,p(s_{t+1}\mid s_t,a_t)
\qquad\qquad
\max_\theta\ \mathbb{E}_{p_{\pi_\theta}}\left[\sum_t r(s_t,a_t)\right]$$

\begin{center}
\begin{tabular}{@{}p{3.2cm}p{10.3cm}@{}}
\toprule
Symbol & What it means \\
\midrule
$\mathcal{D}: \{(s,a,s',r)\}$ & The offline dataset: recorded (state, action, next
state, reward) tuples. \\
$\pi_\beta$ & The unknown ``behavior policy'' that actually generated the dataset
(possibly a mixture of several policies). \\
$p(s_1)$ & The probability distribution over starting states. \\
$p(s_{t+1}\mid s_t,a_t)$ & The environment's own dynamics: how the world responds to an
action (not something we control or choose). \\
$p_{\pi_\theta}$ & The distribution over trajectories that would occur if we actually ran
our trained policy $\pi_\theta$ in the environment. \\
$r(s_t,a_t)$ & The reward received at each step. \\
\bottomrule
\end{tabular}
\end{center}

\paragraph{What the equations are saying.}
The first equation just says: a recorded trajectory's probability is built from the
starting-state distribution, the behavior policy's action choices, and the environment's
own dynamics --- multiplied together over every time step. The second equation is the
important one: our objective is expected total reward under $\pi_\theta$ (the policy we
are training), \emph{not} under $\pi_\beta$ (the policy that collected the data). This
mismatch between ``what generated the data'' and ``what we're evaluating'' is the whole
source of difficulty in offline RL.

\paragraph{Aspects of this concept.}
\begin{itemize}
\item \textbf{What problem it defines:} formalizes offline RL's goal precisely so we can
reason about whether an algorithm is even trying to solve the right problem.
\item \textbf{Key subtlety:} because the objective's expectation is under $\pi_\theta$
while the only data we have came from $\pi_\beta$, any algorithm that is not careful can
end up trusting parts of $\pi_\theta$'s behavior that were never actually observed.
\item \textbf{What if this were done differently?} If we instead maximized the expected
reward under $p_{\pi_\beta}$ (the data's own distribution), we would simply be
\emph{evaluating} the behavior policy that already exists, not improving on it. The
entire value of offline RL comes from doing better than $\pi_\beta$ while only ever
looking at $\pi_\beta$'s data.
\end{itemize}

\section{Slide 16: Offline RL versus Imitation Learning --- Stitching}

\paragraph{What is the concept?}
This slide makes the central case for offline RL over imitation learning: a good offline
RL algorithm can combine (``stitch together'') good pieces from \emph{different}
trajectories into a single new trajectory that is better than any trajectory actually
demonstrated. Plain imitation learning can never do this.

\paragraph{Explanation.}
Imagine two separate recorded trips through a maze of connected states. One trip
successfully reaches the goal ($s_9$) after passing through $s_7$; a different, separate
trip passes through $s_3\to s_7$ but then goes somewhere unhelpful and ends badly. Even
though no single recorded trip goes cleanly from the start all the way to the goal via
$s_3\to s_7\to s_9$, a good offline RL algorithm can realize that ``$s_7$ leads somewhere
great'' (learned from the first trip) and connect that fact to ``$s_3\to s_7$ is a
transition that exists in the data'' (from the second trip) --- gluing two half-good
trips into one fully good path.

\paragraph{The key claim on the slide.}
\emph{Imitation methods cannot outperform the policy that collected the data. Offline
RL can leverage reward information to outperform the behavior policy. Good offline RL
methods can stitch together good behaviors.}

\paragraph{Mechanism of stitching (why it works, step by step).}
\begin{enumerate}
\item From the trajectory that reaches the goal via $s_7\to s_9$ (with reward $r(s_9)=+1$),
temporal-difference (TD) learning assigns a high learned value to $s_7$, because that
value backs up from the good reward at $s_9$.
\item A \emph{different} trajectory contains the transition $s_3\to s_7$. Its TD target
for that step is $r + \gamma V(s_7)$ --- it uses the \emph{learned} value of $s_7$, not
whatever return that second trajectory itself happened to observe afterward.
\item Because $V(s_7)$ is high (learned from the \emph{first} trajectory), the
transition $s_3\to s_7$ gets credited with a high value, even though the trajectory it
was recorded in did not end well.
\item This is only possible because of \textbf{bootstrapping} (TD learning): the value
used in the target comes from the learned value function, not from what actually
happened afterward in that particular trajectory.
\end{enumerate}

\paragraph{Does it matter whether we use TD or Monte Carlo? Yes.}
Monte Carlo (MC) value estimation would instead credit the $s_3\to s_7$ transition with
the \emph{actual observed return} of the trajectory it came from (which ended badly) ---
it has no mechanism for borrowing the good outcome discovered in a completely different
trajectory. So stitching \emph{depends specifically on using TD}, not MC.

\paragraph{Aspects of this concept.}
\begin{itemize}
\item \textbf{Advantage of offline RL over plain imitation learning:} can exceed the
performance of the data-collecting policy by recombining good sub-behaviors that were
never demonstrated together in a single trajectory.
\item \textbf{Why this is the deepest idea in this section (justify-the-design):}
stitching is not a side benefit of TD learning --- it is the \emph{direct consequence} of
using a learned, bootstrapped value target instead of the trajectory's own observed
return.
\item \textbf{What if this were done differently?} If value estimation used Monte Carlo
returns instead of TD bootstrapping, stitching would not happen --- this is exactly why
AWR (Slide 21), which uses Monte Carlo, cannot outperform the behavior policy the way
AWAC or IQL (which use TD) potentially can.
\end{itemize}

\section{Slides 17--18: Why Not Just Use Ordinary Off-Policy Algorithms?}

\paragraph{What is the concept?}
These two slides show what happens if we take an ordinary off-policy actor-critic method
(like SAC, designed for \emph{online} RL with a replay buffer) and simply apply it to a
fixed, static, offline dataset. The result is a specific, serious failure: the learned
Q-function becomes badly overestimated on actions that were never in the data.

\paragraph{Explanation.}
The usual actor-critic recipe fits a Q-function (``how good is this action here?'') and
then improves the policy to pick whatever action the Q-function currently says is best.
When training online, if the Q-function is wrong about some action, the agent can simply
\emph{try it} and find out the truth, correcting the mistake. But offline, there is no
real environment to check against. If the Q-function is (wrongly) very optimistic about
some action that never appears in the dataset, the policy will happily chase that
mistake, and nothing will ever correct it.

\paragraph{The equations.}
$$\min_\phi \sum_{(s,a,s')\sim\mathcal{D}} \left\| \hat Q_\phi^{\pi_\theta}(s,a) -
\Big(r(s,a) + \gamma\,\mathbb{E}_{a'\sim\pi_\theta(\cdot\mid s')}\big[\hat Q_\phi^{\pi_\theta}(s',a')\big]\Big)\right\|^2$$
$$\nabla_\theta J(\theta) \approx \frac{1}{N}\sum_i \nabla_\theta \log \pi_\theta(a_i^{\pi^*}\mid s_i)\,
\hat Q_\phi^{\pi_\theta}(s_i, a_i^{\pi^*}), \qquad a_i^{\pi^*}\sim \pi_\theta(a\mid s_i)$$

\begin{center}
\begin{tabular}{@{}p{3.2cm}p{10.3cm}@{}}
\toprule
Symbol & What it means \\
\midrule
$\hat Q_\phi^{\pi_\theta}(s,a)$ & The learned critic: an estimate of how good action $a$
is in state $s$, for the current policy $\pi_\theta$. \\
$a' \sim \pi_\theta(\cdot\mid s')$ & A \emph{freshly sampled} next action from the
\emph{current} policy --- not an action read from the dataset. \\
$a_i^{\pi^*} \sim \pi_\theta(a\mid s_i)$ & A freshly sampled action from the current
policy, used to improve the actor by climbing the critic. \\
Out-of-distribution (OOD) action & An action that the dataset never actually contains
for a given state. \\
Data support & The range of actions the dataset actually contains for a given state. \\
\bottomrule
\end{tabular}
\end{center}

\paragraph{What the equations are saying.}
The critic update tries to make $\hat Q$ consistent with the reward plus the value of
whatever action the \emph{current} policy would pick next (not the action actually
recorded in the data). The actor update then pushes the policy toward whatever action
currently maximizes $\hat Q$. Both steps repeatedly ask the critic about actions the
current policy \emph{chooses}, which may well not be actions ever seen in the dataset.

\paragraph{Mechanism of the failure (the OOD chain --- memorize this in order).}
\begin{enumerate}
\item The critic's bootstrapped target requires $\mathbb{E}_{a'\sim\pi_\theta(\cdot\mid
s')}[\hat Q(s',a')]$, and $a'$ need not appear anywhere in $\mathcal{D}$.
\item $\hat Q$ is unreliable --- essentially untrained noise --- on such out-of-distribution
(OOD) actions, since nothing in the data constrains its value there.
\item Because the actor update explicitly maximizes $\hat Q$, it actively seeks out
exactly the actions where $\hat Q$ happens to be (wrongly) high.
\item Those inflated, incorrect values then feed into the \emph{next} round's TD target,
so Q-values across the board become substantially overestimated.
\item The learned policy drifts further and further from the real behavior policy, with
no real environment interaction available to ever catch or correct the mistake. Online
RL escapes this loop by actually trying the action and observing the true reward;
offline, the loop stays closed forever. This is exactly why it is the \emph{offline}
setting specifically that breaks this algorithm.
\end{enumerate}

\paragraph{Aspects of this concept.}
\begin{itemize}
\item \textbf{Advantage of this recipe (in its native, online setting):} efficient reuse
of a replay buffer while still being able to correct mistakes through live interaction.
\item \textbf{Disadvantage (offline specifically):} catastrophic, self-reinforcing
Q-value overestimation.
\item \textbf{When not to use this recipe as written:} whenever training purely offline
on a fixed, static dataset.
\item \textbf{When it is fine to use:} online settings, where the agent can actually try
the action it is curious about and observe the real reward.
\item \textbf{What if this were done differently?} If $a'$ were instead sampled from the
dataset itself ($a'\sim\mathcal{D}$) rather than from $\pi_\theta$, the OOD query
disappears, but now the critic is evaluating something closer to $Q^{\pi_\beta}$ (the
old behavior policy) rather than the current policy $\pi_\theta$ --- trading one problem
(OOD overestimation) for another (evaluating the wrong policy). This exact trade-off
reappears explicitly as a noted ``alternative'' in AWAC (Slide 23).
\end{itemize}

\section{Slide 19: Offline RL --- A Baseline: Filtered Behavior Cloning}

\paragraph{What is the concept?}
A very simple way to use reward information without touching Q-functions or value
functions at all: rank whole trajectories by their total reward, keep only the good
ones, and run ordinary behavior cloning on what remains.

\paragraph{Explanation.}
Score every recorded trip by how much total reward it earned. Throw away the low-scoring
trips. Pretend the remaining, better trips are the only dataset that ever existed, and
run plain behavior cloning (Slide 3/6) on it.

\paragraph{The equations.}
$$r(\tau) = \sum_{(s_t,a_t)\in\tau} r(s_t,a_t) \qquad
\widetilde{\mathcal{D}} : \{\tau \mid r(\tau) > \eta\} \qquad
\max_\theta \sum_{(s,a)\in\widetilde{\mathcal{D}}} \log \pi_\theta(a\mid s)$$

\begin{center}
\begin{tabular}{@{}p{3.2cm}p{10.3cm}@{}}
\toprule
Symbol & What it means \\
\midrule
$r(\tau)$ & The total reward earned across an entire trajectory $\tau$. \\
$\eta$ & A threshold return; trajectories must beat this to be kept. \\
$\widetilde{\mathcal{D}}$ & The filtered dataset: only the trajectories with
$r(\tau) > \eta$. \\
$\max_\theta \sum \log \pi_\theta(a\mid s)$ & Ordinary maximum-likelihood behavior
cloning, applied only to $\widetilde{\mathcal{D}}$. \\
\bottomrule
\end{tabular}
\end{center}

\paragraph{What the equations are saying.}
First, add up the reward across each whole trip. Second, keep only the trips whose total
beats some cutoff $\eta$ (e.g., the top $k\%$ of trips). Third, forget about rewards
entirely and just imitate whatever is left, exactly like ordinary behavior cloning.

\paragraph{Aspects of this concept.}
\begin{itemize}
\item \textbf{Advantage:} extremely simple; no value function, no Q-learning, and
therefore no out-of-distribution query risk at all --- it fully reuses the plain
behavior-cloning machinery.
\item \textbf{Disadvantage:} filtering happens at the level of the \emph{whole
trajectory}. If a trajectory clears the threshold, \emph{every} $(s,a)$ pair inside it
survives into $\widetilde{\mathcal{D}}$ --- including any individual bad or mediocre
actions that happened to occur inside an overall good trip. Also, if only one kept
trajectory ever visits a particular state, the learned action there is entirely
determined by that one trajectory, with no evidence that it was actually the best
possible action at that state.
\item \textbf{Problem it solves:} a low-risk way to make some use of reward
information, without any of the machinery (or the OOD danger) of Q-learning.
\item \textbf{Problem it cannot solve:} it cannot stitch trajectories together (Slide
16), and it cannot tell a genuinely good action apart from a merely mediocre action that
happened to ride along inside an otherwise-good trajectory.
\item \textbf{When to use it:} as a quick, low-risk baseline when return-labeled
trajectories are available.
\item \textbf{When not to use it:} when very few trajectories clear the threshold, or
when good actions are scattered as isolated good moments inside otherwise mediocre
trajectories (that signal is lost, since only whole trajectories are judged).
\item \textbf{What if this were done differently?} If filtering were applied at the
level of \emph{individual} $(s,a)$ pairs (e.g., by a per-step advantage estimate)
instead of whole trajectories, we would move toward weighted imitation learning (Slide
20) rather than filtered behavior cloning.
\end{itemize}

\section{Slide 20: Weighted Imitation Learning}

\paragraph{What is the concept?}
Instead of a hard yes/no cutoff (filtered BC), give every individual $(s,a)$ pair a
\emph{soft} weight based on how much better than average that specific action was ---
its \textbf{advantage}.

\paragraph{Explanation.}
Recall the advantage function: $A_\pi(s,a) = Q_\pi(s,a) - V_\pi(s)$, i.e., how much
better action $a$ is compared to the average action in state $s$. Weighted imitation
learning uses this number to decide how hard to imitate each recorded action: actions
with a big positive advantage get imitated strongly (a big weight); actions with an
advantage near zero or negative barely influence training at all, but are not thrown
away completely the way filtered BC would throw away an entire trajectory.

\paragraph{The equations.}
$$A_\pi(s,a) = Q_\pi(s,a) - V_\pi(s) \qquad\qquad
\theta \leftarrow \arg\max_\theta\ \mathbb{E}_{s,a\sim\mathcal{D}}\left[\log \pi_\theta(a\mid s)\,
\exp\big(A(s,a)\big)\right]$$

\begin{center}
\begin{tabular}{@{}p{3.2cm}p{10.3cm}@{}}
\toprule
Symbol & What it means \\
\midrule
$Q_\pi(s,a)$ & Expected total future reward if you take action $a$ in state $s$, then
follow $\pi$ afterward. \\
$V_\pi(s)$ & Expected total future reward from state $s$, averaged over the actions
$\pi$ would itself choose. \\
$A_\pi(s,a)$ & How much better (or worse) action $a$ is than the average action in
state $s$. \\
$\exp(A(s,a))$ & The soft weight given to this $(s,a)$ pair during imitation: always
positive, growing quickly for large positive advantages. \\
\bottomrule
\end{tabular}
\end{center}

\paragraph{What the equations are saying.}
The first equation defines ``how good was this action, relative to average.'' The
second equation says: do ordinary maximum-likelihood imitation (as in Slide 6), but
multiply each example's contribution by $\exp(A(s,a))$ --- so better-than-average
actions get imitated much more strongly than merely average or below-average ones.

\paragraph{Aspects of this concept.}
\begin{itemize}
\item \textbf{Advantage over filtered BC:} smoother and more fine-grained --- it uses
information from \emph{every} sample rather than an all-or-nothing cutoff, and it
discriminates at the level of individual $(s,a)$ pairs rather than whole trajectories.
\item \textbf{Disadvantage:} needs an advantage estimate $A(s,a)$, which itself must be
learned from a $V$ and/or $Q$ function --- introducing its own source of estimation
error. (Slides 21--26 are all about \emph{how} to estimate this advantage well.)
\item \textbf{Problem it solves:} gives a smoother, more data-efficient way to bias
imitation toward better actions than a hard trajectory-level filter.
\item \textbf{Problem it does not solve by itself:} it does not specify how to estimate
$A(s,a)$ --- that choice (Monte Carlo vs.\ TD, on which policy) is exactly what
distinguishes AWR, AWAC, and IQL from each other.
\item \textbf{When to use it:} whenever a reasonable $Q$/$V$ estimate is available and
you want a policy-gradient-free way to bias imitation toward good behavior.
\item \textbf{What if this were done differently?} If the weight were simply $A(s,a)$
instead of $\exp(A(s,a))$, negative advantages would produce \emph{negative} weights,
which is not a valid weighting for a log-likelihood (weights in weighted MLE must stay
positive). The exponential also makes large advantages dominate multiplicatively,
sharpening the imitation toward the very best actions.
\end{itemize}

\section{Slide 21: Advantage-Weighted Regression (AWR)}

\paragraph{What is the concept?}
A complete, two-step algorithm that fills in weighted imitation learning (Slide 20) with
a concrete way to estimate the advantage: fit a value function using plain Monte Carlo
regression, then run weighted imitation with a temperature knob $\alpha$.

\paragraph{Explanation.}
Step 1: for every state, look at the total reward that was \emph{actually observed}
afterward in the data, and train $\hat V$ to regress toward that number directly --- an
ordinary supervised-learning target, no bootstrapping. Step 2: run weighted imitation
exactly as in Slide 20, where the ``how good was this action'' score is (the observed
future return) minus (this state's estimated average value), divided by a knob $\alpha$
before being exponentiated, so the weights do not explode into astronomically large
numbers.

\paragraph{The equations.}
$$\min_\phi \sum_{s_t\sim\mathcal{D}} \left\| \hat V_\phi^{\pi_\beta}(s_t) -
\sum_{t'=t}^{T} r(s_{t'},a_{t'})\right\|^2$$
$$\max_\theta\ \mathbb{E}_{(s_t,a_t)\sim\mathcal{D}}\left[\log \pi_\theta(a_t\mid s_t)\,
\exp\left(\frac{1}{\alpha}\left(\sum_{t'=t}^{T} r(s_{t'},a_{t'}) -
\hat V_\phi^{\pi_\beta}(s_t)\right)\right)\right]$$

\begin{center}
\begin{tabular}{@{}p{3.2cm}p{10.3cm}@{}}
\toprule
Symbol & What it means \\
\midrule
$\hat V_\phi^{\pi_\beta}(s_t)$ & The estimated value of state $s_t$ for the
behavior policy $\pi_\beta$ that produced the data (not for $\pi_\theta$). \\
$\sum_{t'=t}^T r(s_{t'},a_{t'})$ & The actual total reward observed from time $t$ to the
end of that trajectory (the Monte Carlo return). \\
$\alpha$ & A temperature hyperparameter that keeps the exponential weight from exploding. \\
\bottomrule
\end{tabular}
\end{center}

\paragraph{What the equations are saying.}
The first equation fits $\hat V$ using ordinary regression toward the observed return
--- exactly like fitting any supervised model, with no learned bootstrap target
involved. The second equation is weighted imitation (Slide 20), where the advantage is
(observed return $-$ estimated value), scaled by $1/\alpha$.

\paragraph{Understanding $\alpha$.}
Small $\alpha$ makes the weighting sharper --- close to only imitating the very best
actions (similar to filtered BC). Large $\alpha$ flattens the weighting --- close to
imitating everything roughly equally, i.e., close to plain, unweighted BC.

\paragraph{Worked example: the deterministic-behavior-policy edge case.}
If the behavior policy $\pi_\beta$ always takes exactly the same single action in a
given state, then the observed return from that state already \emph{equals}
$\hat V^{\pi_\beta}(s)$ for every visit, so the advantage is $0$ for every sample, every
weight is $\exp(0)=1$, and AWR degenerates to exactly plain, unweighted behavior cloning.

\paragraph{Aspects of this concept.}
\begin{itemize}
\item \textbf{Advantage:} very simple value-fitting step (ordinary regression, no
bootstrapping instability); never queries or trains on any out-of-distribution action,
because it only ever uses returns actually observed following data actions.
\item \textbf{Disadvantage 1:} Monte Carlo returns are noisy --- a single trajectory's
full return is a high-variance sample of the true expected value.
\item \textbf{Disadvantage 2:} $\hat V^{\pi_\beta}$ is an estimate for the (typically
weaker) behavior policy $\pi_\beta$, not for the policy $\pi_\theta$ we are actually
trying to improve --- so the weighting is somewhat mismatched with the policy we hope to
end up with.
\item \textbf{The interpretation trap:} a high AWR weight on an action that appears only
once or twice in the dataset should not be trusted --- with so little data, the estimate
of $Q$ (via the observed return) for that action is itself high-variance, and
exponentiating amplifies that noise rather than correcting it. AWR trusts its estimates
blindly; it has no explicit notion of ``how confident should I be here.''
\item \textbf{When to use it:} when a simple, stable, OOD-safe offline algorithm is
wanted and Monte Carlo's variance and behavior-policy mismatch are acceptable.
\item \textbf{What if this were done differently?} If $\hat V$ were instead fit with
TD (bootstrapping) rather than Monte Carlo, and targeted at $\pi_\theta$ rather than
$\pi_\beta$, we would arrive at AWAC (Slide 23) --- lower variance and on-policy-relevant
values, but at the cost of reopening the OOD-query problem.
\end{itemize}

\section{Slide 22: Monte Carlo Estimation --- A Trade-off Summary}

\paragraph{What is the concept?}
A short pros-and-cons summary of using Monte Carlo returns (as AWR does) to estimate
value.

\begin{center}
\begin{tabular}{@{}p{1.2cm}p{12.3cm}@{}}
\toprule
& \\
\midrule
$+$ & Simple: it is ordinary supervised regression toward an observed number. \\
$+$ & Completely avoids querying or training on any out-of-distribution action, since it
only ever uses returns that were literally observed after real, data-recorded actions. \\
$-$ & Monte Carlo estimation is noisy (high variance): a single trajectory's return is
only one sample of a random outcome. \\
$-$ & $\hat A^{\pi_\beta}$ is an advantage estimate for the (weaker) behavior policy,
not for the (better) target policy $\pi_\theta$ we actually want to end up with. \\
\bottomrule
\end{tabular}
\end{center}

\paragraph{Aspects of this concept.}
This slide is the bridge between AWR and AWAC. The two weaknesses listed here --- noisy
estimates and being tied to the wrong (weaker) policy --- are exactly what motivates
replacing Monte Carlo with TD bootstrapping targeted at $\pi_\theta$, which is the whole
idea of the next slide (AWAC). The cost of doing so is reopening the out-of-distribution
problem that Monte Carlo estimation structurally avoided (Slides 17--18) --- this is a
genuine trade-off, not a strict improvement.

\section{Slide 23: Advantage-Weighted Actor-Critic (AWAC)}

\paragraph{What is the concept?}
AWAC uses the exact same weighted-imitation recipe as AWR (Slide 20/21), but replaces
the Monte Carlo value estimate with a TD (bootstrapped) estimate of $Q^{\pi_\theta}$ ---
the value \emph{for the current policy}, not the old behavior policy.

\paragraph{Explanation.}
Instead of waiting to see a whole trajectory's actual total reward (as AWR does), AWAC
learns a Q-function the usual actor-critic way: bootstrap off of a one-step reward plus
the estimated value of whatever the \emph{current} policy would do next. This tends to be
less noisy than a full Monte Carlo return, and (unlike AWR) it is actually an estimate
\emph{for the policy we are improving}, not the old data-collecting policy.

\paragraph{The equations.}
$$\min_\phi\ \mathbb{E}_{(s,a,r,s')\sim\mathcal{D}}\left[\left(\hat Q_\phi^{\pi_\theta}(s,a) -
\Big(r + \gamma\,\mathbb{E}_{a'\sim\pi_\theta(\cdot\mid s')}\big[\hat Q_\phi^{\pi_\theta}(s',a')\big]\Big)\right)^2\right]$$
$$\theta \leftarrow \arg\max_\theta\ \mathbb{E}_{s,a\sim\mathcal{D}}\left[\log \pi_\theta(a\mid s)\,
\exp\big(\hat A^{\pi_\theta}(s,a)\big)\right], \qquad
\hat A^{\pi_\theta}(s,a) = \hat Q_\phi^{\pi_\theta}(s,a) -
\mathbb{E}_{\bar a\sim\pi_\theta(\cdot\mid s)}\left[\hat Q_\phi^{\pi_\theta}(s,\bar a)\right]$$

\begin{center}
\begin{tabular}{@{}p{3.2cm}p{10.3cm}@{}}
\toprule
Symbol & What it means \\
\midrule
$\hat Q_\phi^{\pi_\theta}$ & The critic, estimated via TD (one-step bootstrapping), for
the current policy $\pi_\theta$ (not for the behavior policy $\pi_\beta$ that produced
the data). \\
$a'\sim\pi_\theta(\cdot\mid s')$ & A freshly sampled next action from the current policy
--- this is the out-of-distribution query, noted explicitly in red on the slide. \\
An alternative noted on the slide & $a' \sim \mathcal{D}$: sample the next action from
the dataset instead, to avoid the OOD query (at a cost, discussed below). \\
$\bar a \sim \pi_\theta(\cdot\mid s)$ & A freshly sampled action used only to compute
$V(s)$ as the average $Q$ over the current policy's own choices. \\
\bottomrule
\end{tabular}
\end{center}

\paragraph{What the equations are saying.}
The critic equation is a standard actor-critic TD update, except the ``next action''
$a'$ used in the bootstrap target comes from the \emph{current} policy $\pi_\theta$, not
from the dataset. The policy equation is exactly weighted imitation learning (Slide 20),
with the advantage now coming from this TD-based, current-policy critic instead of
Monte Carlo.

\paragraph{Aspects of this concept.}
\begin{itemize}
\item \textbf{Advantage:} because the critic is $Q^{\pi_\theta}$ (current policy) rather
than $Q^{\pi_\beta}$ (behavior policy), and because TD estimates are typically
lower-variance than Monte Carlo returns, AWAC's advantage estimates are both more
relevant and less noisy than AWR's --- allowing it to potentially improve beyond the
behavior policy in a way AWR structurally cannot. It also still only ever trains the
actor's log-likelihood term on actions actually present in the dataset ($a\in\mathcal{D}$).
\item \textbf{Disadvantage:} the critic's bootstrap target samples
$a'\sim\pi_\theta(\cdot\mid s')$ --- exactly the out-of-distribution query move from
Slides 17--18. AWAC therefore reopens the Q-overestimation problem that AWR completely
avoided. In a fixed offline dataset with no way to check against reality, this is
genuinely risky.
\item \textbf{Trade-off table (the three offline algorithms so far):}
\begin{center}
\begin{tabular}{@{}p{2cm}p{3.3cm}p{2.3cm}p{6.4cm}@{}}
\toprule
& Value estimate & Touches OOD actions? & Weakness \\
\midrule
AWR & Monte Carlo $\hat V^{\pi_\beta}$ & No & Monte Carlo is noisy; the advantage is for
the weaker behavior policy, not $\pi_\theta$. \\
AWAC & TD $\hat Q^{\pi_\theta}$, $a'\sim\pi_\theta$ & Yes (in the critic) & Reintroduces
OOD queries and their overestimation risk. \\
\bottomrule
\end{tabular}
\end{center}
(IQL, the third row of this table, is completed in Slide 26.)
\item \textbf{What if this were done differently?} The slide's own noted alternative ---
sampling $a'\sim\mathcal{D}$ instead of $a'\sim\pi_\theta(\cdot\mid s')$ --- removes the
OOD query, but then the critic is again evaluating something closer to $Q^{\pi_\beta}$,
undoing exactly the advantage AWAC was trying to gain over AWR. This is the same
underlying trade-off already seen in Slides 17--18.
\end{itemize}

\section{Slide 24: Expectile Regression --- Motivation}

\paragraph{What is the concept?}
This slide poses the question that motivates the next two slides: instead of estimating
the \emph{mean} of a random variable (which ordinary squared-error loss always gives
you), can we estimate something closer to a high percentile --- an ``optimistic but
still data-grounded'' summary --- while still only ever looking at actions that actually
appear in the data?

\paragraph{Explanation.}
Picture, for a fixed state $s'$, a histogram of $Q(s',a)$ values for all the different
actions $a$ that the behavior policy actually tried there. Ordinary $\ell_2$ (squared
error) regression will always find you the \emph{average height} of this histogram. But
what we would really like, to avoid ever needing to query an action outside the data, is
something like ``the value of the \emph{best} action that was actually tried'' ---
closer to a high percentile of the histogram, computed \emph{only} using the actions
that are genuinely there.

\paragraph{Aspects of this concept.}
\begin{itemize}
\item \textbf{Problem it is trying to solve:} get an estimate of $V(s)$ that behaves
like ``the best of what was tried'' rather than ``the average of what was tried,''
without ever plugging in an action that was never tried (i.e., without any
$\max_{a'}Q(s,a')$ over all conceivable actions).
\item \textbf{Why ordinary $\ell_2$ loss cannot do this:} squared-error loss is
minimized exactly at the mean --- it has no built-in way to lean toward the high or low
end of a distribution.
\item \textbf{What problem it gives rise to:} the need for an asymmetric loss function
that can be pushed toward a higher or lower ``expectile'' --- this is exactly the
expectile loss defined on the next slide.
\end{itemize}

\section{Slide 25: The Expectile Loss}

\paragraph{What is the concept?}
A generalization of ordinary squared-error loss with one extra knob, $\lambda$, that
lets the fitted value lean toward the upper or lower part of a distribution instead of
always sitting exactly at the mean.

\paragraph{Explanation.}
Ordinary squared error treats being-too-high and being-too-low exactly the same way. The
expectile loss instead makes one direction of error more expensive than the other,
controlled by a single knob $\lambda$ between 0 and 1. Whichever direction is made more
\emph{expensive}, the best-fit value will shift \emph{away} from making that kind of
error --- which pulls the fit toward the \emph{opposite} side of the distribution.

\paragraph{The equation.}
$$\ell_2^\lambda(x) = \begin{cases}(1-\lambda)\,x^2 & \text{if } x<0\\ \lambda\, x^2 &
\text{otherwise}\end{cases}$$

\begin{center}
\begin{tabular}{@{}p{3.2cm}p{10.3cm}@{}}
\toprule
Symbol & What it means \\
\midrule
$x$ & The residual (error): what is being subtracted from what depends on the specific
application (defined explicitly for IQL on the next slide). \\
$\lambda$ & The expectile level, a number between 0 and 1 that decides which direction
of error is penalized more. \\
$\ell_2^\lambda(x)$ & The loss value: an ordinary squared error, but scaled differently
depending on the sign of $x$. \\
\bottomrule
\end{tabular}
\end{center}

\paragraph{What the equation is saying.}
If $x$ is negative, the error is scaled by $(1-\lambda)$; if $x$ is zero or positive, it
is scaled by $\lambda$. At $\lambda=0.5$, both sides are scaled equally and this is
exactly ordinary squared error (which gives back the mean --- $\lambda = 0.5$ is the
``50th expectile,'' analogous to the median). As $\lambda$ moves toward $1$, positive
$x$ becomes very expensive, so the best fit shifts \emph{down} to avoid producing
positive residuals as often (a lower expectile). As $\lambda$ moves toward $0$, negative
$x$ becomes very expensive, so the best fit shifts \emph{up} (a higher, more optimistic
expectile).

\paragraph{Aspects of this concept.}
\begin{itemize}
\item \textbf{What it is:} an asymmetric squared-error loss; a strict generalization of
ordinary mean-squared error, which is the special case $\lambda=0.5$.
\item \textbf{Advantage:} lets us target an ``expectile'' (a percentile-like statistic)
instead of only the mean, while remaining smooth and easy to optimize with gradients ---
unlike quantile regression, which uses an absolute-value loss that is not smooth at
zero.
\item \textbf{Disadvantage:} an expectile is a related but technically different
statistic from a literal percentile (it is defined by this asymmetric-loss minimization,
not by sorting the data), so the intuition ``$\lambda=0.9$ means the 90th percentile'' is
only approximate; it also introduces a new hyperparameter, $\lambda$, that must be
chosen.
\item \textbf{Problem it solves:} gives a tunable knob for whether a fitted value leans
optimistic (upper expectile) or pessimistic (lower expectile).
\item \textbf{Problem it does not solve by itself:} this slide does not say what $x$
actually represents in a specific algorithm --- that choice is made on the very next
slide, for IQL.
\item \textbf{When to use it:} whenever you want a value estimate that reflects
near-best (or near-worst) outcomes rather than the average, using only the data you
actually have.
\item \textbf{What if this were done differently?} Using an absolute-value (quantile)
loss instead of a squared-error-based one would give literal quantiles rather than
expectiles, at the cost of losing the smooth, easy-to-optimize gradient that squared
error provides.
\end{itemize}

\section{Slide 26: Implicit Q-Learning (IQL)}

\paragraph{What is the concept?}
A full, three-step offline RL algorithm that uses expectile regression to estimate an
optimistic value function \emph{without ever evaluating $Q$ on an action outside the
dataset} --- completing the trade-off table started in Slide 23.

\paragraph{Note on the slide.}
Compare this slide's formula for $V$, which uses $x = V(s) - \hat Q(s,a)$ and instructs
``using small $\lambda<0.5$,'' against the expectile loss on Slide 25, which used
$x<0$ to mean the $(1-\lambda)$-weighted branch. Because $x = V - Q$ here, $x<0$ means
$V$ is \emph{below} $Q$ (i.e., $V$ underestimates $Q$). With small $\lambda$, that
branch is weighted by $(1-\lambda)$, which is close to $1$ --- i.e., underestimating $Q$
is made \emph{expensive}. So the best-fit $V$ shifts \emph{up}, toward the upper
expectile, exactly as intended. This is self-consistent with the slide's own
instruction. It is worth knowing, however, that this labeling is the mirror image of how
the original IQL paper writes the same loss (using $x = Q - V$ and a $\lambda$ close to
$1$ to mean the upper expectile) --- both conventions produce the identical intended
behavior (an optimistic $V$), they simply place the label ``small $\lambda$'' versus
``large $\lambda$'' on opposite sides. On the exam, state which convention you are using
in one sentence, then answer using the mechanism (``$V$ shifts above the mean''), not
just the raw label.

\paragraph{Explanation of the full pipeline.}
\begin{enumerate}
\item Fit $\hat V(s)$ to be an optimistic-but-honest summary of $\hat Q$ at each state,
using expectile regression, restricted entirely to the $(s,a)$ pairs that are actually in
the dataset. This behaves like a soft ``best of what was tried'' summary.
\item Fit $\hat Q(s,a)$ with an ordinary squared-error TD update, using $\hat V(s')$ ---
not a max or expectation over next actions --- as the bootstrap target. This is the
crucial trick: it never needs to evaluate $Q$ at any next action at all.
\item Extract an actual usable policy using the same weighted-imitation recipe as AWR
and AWAC.
\end{enumerate}

\paragraph{The equations.}
$$\hat V(s) \leftarrow \arg\min_V\ \mathbb{E}_{(s,a)\sim\mathcal{D}}\left[\ell_2^\lambda
\big(V(s) - \hat Q(s,a)\big)\right] \quad\text{(using small } \lambda<0.5\text{)}$$
$$\hat Q(s,a) \leftarrow \arg\min_Q\ \mathbb{E}_{(s,a,s')\sim\mathcal{D}}\left[\left(Q(s,a) -
\big(r + \gamma \hat V(s')\big)\right)^2\right]$$
$$\hat\pi \leftarrow \arg\max_\pi\ \mathbb{E}_{s,a\sim\mathcal{D}}\left[\log \pi(a\mid s)\,
\exp\left(\frac{1}{\alpha}\big(\hat Q(s,a) - \hat V(s)\big)\right)\right]$$

\begin{center}
\begin{tabular}{@{}p{3.2cm}p{10.3cm}@{}}
\toprule
Symbol & What it means \\
\midrule
$\hat V(s)$ & The optimistic value estimate, fit via expectile regression using only
dataset $(s,a)$ pairs. \\
$\hat Q(s,a)$ & The Q-function, fit with plain TD (squared-error) regression, bootstrapped
off $\hat V(s')$ instead of any action-dependent next-step quantity. \\
$\lambda$ & The expectile level for the $V$-step; small (as instructed) pushes $\hat V$
toward the upper expectile of $\hat Q$ over dataset actions. \\
$\alpha$ & The temperature for the final weighted-imitation policy-extraction step
(same role as in AWR/AWAC). \\
\bottomrule
\end{tabular}
\end{center}

\paragraph{Why IQL never touches OOD actions (the key justify-the-design point).}
Standard Q-learning's bootstrap target needs $\max_{a'} Q(s',a')$ or
$\mathbb{E}_{a'\sim\pi_\theta}[Q(s',a')]$ --- both of which require evaluating $Q$ at
actions that might never appear in the dataset. IQL replaces that entire
max/expectation with $\hat V(s')$, and $\hat V$ is fit using an expectile loss computed
\emph{only} over the $(s,a)$ pairs that are actually present in $\mathcal{D}$. The
``soft max-like,'' optimistic behavior comes entirely from re-weighting
\emph{existing} data through the asymmetric loss, never from plugging in an unseen
action. This is why IQL avoids out-of-distribution queries \emph{by construction}, not
by luck.

\paragraph{Aspects of this concept.}
\begin{itemize}
\item \textbf{Advantage:} gets the low-variance, current-policy-relevant benefits of TD
bootstrapping (like AWAC) while completely avoiding the OOD-query problem that AWAC
reintroduces.
\item \textbf{Disadvantage:} introduces an extra hyperparameter, $\lambda$ (the
expectile level), which must be tuned; pushed too aggressively toward the extreme
expectile, it can become overconfident about state-action pairs that are only sparsely
represented in the data, since expectile regression from very few noisy samples is
itself imprecise.
\item \textbf{Completed trade-off table:}
\begin{center}
\begin{tabular}{@{}p{2cm}p{3.6cm}p{2.3cm}p{6.1cm}@{}}
\toprule
& Value estimate & Touches OOD? & Weakness \\
\midrule
AWR & Monte Carlo $\hat V^{\pi_\beta}$ & No & Noisy; advantage is for the weaker behavior
policy, not $\pi_\theta$. \\
AWAC & TD $\hat Q^{\pi_\theta}$, $a'\sim\pi_\theta$ & Yes (in the critic) & Reintroduces
OOD queries. \\
IQL & TD with expectile $\hat V$ & No & Extra hyperparameter $\lambda$; conservative,
can be imprecise on sparsely-represented actions. \\
\bottomrule
\end{tabular}
\end{center}
\item \textbf{When to use it:} offline settings where avoiding OOD overestimation
matters most, and TD's efficiency advantage over Monte Carlo is desired.
\item \textbf{When not to use it:} very simple datasets from an essentially
deterministic behavior policy, where plain BC or AWR is already sufficient and IQL's
extra machinery adds complexity without benefit.
\item \textbf{What if this were done differently?} If the $V$-step used plain
$\ell_2$ regression ($\lambda=0.5$) instead of an asymmetric expectile loss, $\hat V$
would just be the ordinary \emph{mean} of $\hat Q$ over the data's actions at that state
--- exactly the plain, non-optimistic estimate that this whole section is trying to move
beyond.
\end{itemize}

\section{Quick Reference: The Progression of Ideas}

\begin{center}
\begin{tabular}{@{}p{3.4cm}p{4.6cm}p{4.6cm}@{}}
\toprule
Starting point & Problem discovered & Fix introduced \\
\midrule
Behavior cloning ($\ell_2$) & Mean-collapse on multimodal experts & Maximum likelihood
with expressive distributions (GMM, autoregressive, diffusion/flow) \\
Expressive behavior cloning & Compounding error / distribution shift & DAgger
(interactive data aggregation) \\
DAgger & Needs an interactive expert, often infeasible & Offline RL: reuse existing
reward-labeled data instead \\
Offline data, naive off-policy actor-critic & Out-of-distribution Q-value
overestimation & Restrict training to actions in the data: filtered BC, then weighted
imitation \\
Filtered BC & Whole-trajectory filtering is too coarse & Weighted imitation with a
per-$(s,a)$ advantage (AWR) \\
AWR (Monte Carlo) & Noisy; advantage is for the weak behavior policy & AWAC (TD,
current-policy advantage) \\
AWAC (TD) & Reopens OOD queries in the critic & IQL (TD with expectile $\hat V$,
no OOD queries) \\
\bottomrule
\end{tabular}
\end{center}

\section{Key Errors and Ambiguities in the Source Slides (Review Before the Exam)}

\begin{enumerate}
\item \textbf{Slide 11, ``Diffusion Policy'':} the procedure shown (linear interpolation
between noise and data, a single learned velocity field, Euler-integrated inference) is
technically \emph{flow matching / rectified flow}, not classic diffusion (DDPM), which
uses many small noising steps and a stochastic reverse process. Answer exam questions
using the slide's own procedure; optionally note the more precise name for extra
precision.
\item \textbf{Slides 25--26, expectile $\lambda$ convention:} Slide 26 defines the
residual as $x = V(s) - \hat Q(s,a)$ and instructs ``using small $\lambda < 0.5$'' to
push $\hat V$ toward the upper expectile (this is internally self-consistent and
verified numerically above). This is the mirror image of the original IQL paper's own
convention, which defines the residual the other way around ($x = Q - V$) and uses
$\lambda$ close to $1$ for the same upper-expectile behavior. State your convention in
one sentence before answering any question that asks about the direction of $\lambda$.
\end{enumerate}

\end{document}